\documentclass
[
a4paper,
german,
openright,                    % Kap.beginn immer rechts! (fkt. nur bei report, nicht bei article)
12pt                          % ersatzweise 12pt, wenn mehr Seiten entstehen sollen
]
{report}
%%%%%%%%%%%%%%%%%%%%%%%%%%%%%%%%%%%%%%%%%%%%%%%%%%%%%%%%%%%%%%%%%%%%%%%%%%%%%%%%%%%%%%%%%%%%%%%%%%%%%%%%%%%%
%Dokumenteneigenschaften

\newcommand{\Author}{Refik Kerimi} 
\newcommand{\Title}{Analyse der Auswirkung von Progressive Web Apps auf bestehende Web Apps}
\newcommand{\Keywords}{Progressive Web App, Manifest.json, Service Workers, Home Banner, Offline}
\newcommand{\Advisor}{DI Norbert Egger BSc}
\newcommand{\Birthdate}{28.03.1983}
\newcommand{\EnrolNum}{1410555043}
\newcommand{\VenueMonthYear}{Salzburg, September 2018}
\newcommand{\VenueDate}{Salzburg, am 1.09.2018}

\newcommand{\frontmatter}{\cleardoublepage\pagenumbering{roman}}

\newcommand{\mainmatter}{\cleardoublepage\pagenumbering{arabic}\setcounter{page}{1}}

%%%%%%%%%%%%%%%%%%%%%%%%%%%%%%%%%%%%%%%%%%%%%%%%%%%%%%%%%%%%%%%%%%%%%%%%%%%%%%%%%%%%%%%%%%%%%%%%%%%%%%%%%%%%
% Layout zusammengestellt von Richard Wanger im WS 2015/2016. Verschiedene Vorlagen wurden bei der Erstellung
% kombiniert und nach dem Masterleitfaden (Stand: Oktober 2015) adaptiert. Verwendung ohne Gewähr!
% Umbau des Layouts nach dem Bachelorleitfaden (Stand: Juni 2018) durch Refik Kerimi SS 2018 

\usepackage{style}		%alle packages und Formatvorlagen befinden sich in dieser Datei


%%%%%%%%%%%%%%%%%%%%%%%%%%%%%%%%%%%%%%%%%%%%%%%%%%%%%%%%%%%%%%%%%%%%%%%%%%%%%%%%%%%%%%%%%%%%%%%%%%%%%%%%%%%%
% ORGANISATORISCHES

\usepackage{amsmath}
\usepackage{amsfonts}
\usepackage{amssymb}


\begin{document}

\begin{titlepage}

\hspace{7cm}

\begin{center}
	{\Large\uppercase\expandafter{\bf BACHELORARBEIT}}\\[0.5ex]
	\vspace{1cm}
	\Large{\bf\large \Title}\\
	\vspace{1.5cm}
	\normalsize durchgeführt am\\
	Studiengang Informationstechnik \& System--Management\\
	an der\\
	Fachhochschule Salzburg GmbH\\
\end{center}

\vspace{2cm}

\begin{center}
	\normalsize vorgelegt von
	\\
	{
		\Large{\bf\large \Author}\\
	}
	\vspace{2cm}
	\includegraphics[width=7cm]{BilderAllgemein/Logo.jpg}\medskip
\end{center}
	


\begin{tabbing}
	\hspace*{3cm}\=\hspace*{4cm}\= \kill
	\> Studiengangsleiter: \> FH-Prof.~DI Dr. Gerhard Jöchtl \\*[0.2cm]
	\> Betreuer: \> \Advisor
\end{tabbing}

\vfill	

\begin{center}
\VenueMonthYear\\
\end{center}
\end{titlepage}
\chapter*{Eidesstattliche Erklärung}
\thispagestyle{plain}
\pagestyle{plain}

Ich versichere an Eides statt, dass ich die vorliegende Bachelorarbeit ohne unzulässige fremde Hilfe und ohne Benutzung anderer als der angegebenen Quellen und Hilfsmittel angefertigt und alle aus ungedruckten Quellen, gedruckter Literatur oder aus dem Internet im Wortlaut oder im wesentlichen Inhalt übernommenen Formulierungen und Konzepte gemäß den Richtlinien wissenschaftlicher Arbeiten zitiert, bzw. mit genauer Quellenangabe kenntlich gemacht habe. Diese Arbeit wurde in gleicher oder ähnlicher Form weder im In- noch im Ausland in irgendeiner Form als Prüfungsarbeit vorgelegt und stimmt mit der durch die Begutachter beurteilten Arbeit überein.
\vspace{3cm}

\VenueDate

\vspace{0.5cm}

\begin{tabular}{p{0.3\textwidth}p{0.32\textwidth}p{0.3\textwidth}}



\parbox[c]{1em}{\includegraphics[width=5cm]{BilderAllgemein/unterschriftRef.jpg}} &  &  \multicolumn{1}{c}{1410555043} \\ \cline{1-1} \cline{3-3}

Refik Kerimi & & Matrikelnummer


\end{tabular}
\chapter*{Allgemeine Informationen}
\thispagestyle{plain}
\pagestyle{plain}
\renewcommand{\footrulewidth}{0.4pt}
\lfoot{\small Refik Kerimi}

\begin{tabular}{p{0.3\textwidth}p{0.65\textwidth}}

Vor- und Zuname: & \Author \\*[0.2cm]
Institution: & Fachhochschule Salzburg GmbH \\*[0.2cm]
Studiengang: & Informationstechnik \& System-Management \\*[0.2cm]
Titel der Bachelorarbeit: & \Title \\*[0.2cm]
Schlagwörter: & \Keywords  \\*[0.2cm]
Betreuer an der FH: & \Advisor

\end{tabular}

\newpage

\section*{\Large\bfseries Kurzfassung}



\section*{\Large\bfseries Abstract}



\chapter*{Danksagung}
\thispagestyle{plain}
\pagestyle{plain}

Danken möchte ich vor allem meinem Betreuer für die Unterstützung bei dieser Bachelorarbeit. 

Besonderer Dank gilt auch meiner Familie und Freunden, die uns während des Studiums in allen Belangen immer unterstützt haben. 



































%%%%%%%%%%%%%%%%%%%%%%%%%%%%%%%%%%%%%%%%%%%%%%%%%%%%%%%%%%%%%%%%%%%%%%%%%%%%%%%%%%%%%%%%%%%%%%%%%%%%%%%%%%%%
%VERZEICHNISSE

\tableofcontents
%\protect\addcontentsline {toc}{chapter} {Inhaltsverzeichnis}

\frontmatter
\renewcommand{\nomname}{Abkürzungsverzeichnis}
\chapter*{Abkürzungsverzeichnis}
\addcontentsline{toc}{chapter}{Abkürzungsverzeichnis} 
\pagestyle{plain}

\begin{acronym}[AUTOSAR]

 \acro {PWA}  				{Progressive Web Applikation}

\end{acronym} 

\listoffigures
\protect \addcontentsline {toc}{chapter}{Abbildungsverzeichnis}

\listoftables
\protect \addcontentsline{toc}{chapter}{Tabellenverzeichnis}

\lstlistoflistings
\protect \addcontentsline{toc}{chapter}{Listingverzeichnis}

%%%%%%%%%%%%%%%%%%%%%%%%%%%%%%%%%%%%%%%%%%%%%%%%%%%%%%%%%%%%%%%%%%%%%%%%%%%%%%%%%%%%%%%%%%%%%%%%%%%%%%%%%%%%
%INHALT

\mainmatter
\chapter{Einleitung}\label{chap:Einleitung}
\thispagestyle{standard}
\pagestyle{standard}
\renewcommand{\footrulewidth}{0.4pt}
\lfoot{\small Refik Kerimi}




\section{Motivation} \label{sub:Motivation}





\newpage
\section{Zielsetzung}\label{sub:Zielsetzung}


\chapter{Grundlagen}
\thispagestyle{standard}
\pagestyle{standard}
\renewcommand{\footrulewidth}{0.4pt}
\lfoot{\small Refik Kerimi}


\section{}

\chapter{Features und Merkmale}\label{chab:FeaturesundMerkmale}
\thispagestyle{standard}
\pagestyle{standard}
\renewcommand{\footrulewidth}{0.4pt}
\lfoot{\small Refik Kerimi}
In diesem Kapitel werden die Komponenten der \acl{PWA} (\acs{PWA}) erklärt. Weiters werden durch die folgenden Tabellen die wichtigsten Punkte zwischen den verschiedenen Technologien  gegenübergestellt.

\section{Aufbau Progressive Web Apps (PWA)}
Die \acs{PWA}s sind keine neuen Technologien, vielmehr sind es verbesserte Stategien, Methoden und APIs wie in Abbildung \ref{fig:Komponenten} zu sehen. 
Sie erleichtern dem User die Benutzung und den Zugriff einer \acs{Web-App} \cite{AlternativePWA}. 

\begin{figure}[h]
	\centering
	\includegraphics[width=14cm]{BilderAllgemein/PWA_Features}\medskip
	\caption{PWA Komponenten}
	\label{fig:Komponenten}
\end{figure}

\section{Unterschiede PWA, Native Applikation und Web-Apps}\label{chap:UnterschiedePWA,NativeApplikationundWeb-Apps}
In den folgenden Tabellen \ref{tab:PwaNvaWaInstallation}, \ref{tab:PwaNvaWaZugriff} und \ref{tab:PwaNvaWaFunktionen}  wird versucht die Features und Merkmale gegenüberzustellen um die Auswahl zu erleichtern. Die Unterschiede in den Punkten Veröffentlichung, Installation, Zugriff und Funktionen werden verglichen.

%Folgende Punkte werden verglichen:
%\begin{itemize}
%    \item  \textbf{Veröffentlichung und Installation}
%	\item  \textbf{Zugriff}
%	\item  \textbf{Funktionen}
%\end{itemize}. 

\begin{table}[h]
\centering

\begin{tabular} {|p{3cm}|p{3.5cm}|p{3.5cm}|p{3.5cm}|}
\hline\multirow{3}{*}
 										&PWA  & Native & Web App	\\ \hline
Veröffentlichung & Es werden verschiedene Entwicklerkonten benötigt Play Store und Apple Store & keine Entwicklerkonten benötigt & keine Entwicklerkonten benötigt\\ \hline

Installation & App muss aus einem der App-Stores downgeloaded werden  & Wird mit einem Klick auf dem Startbildschirm hinzugefügt & keine Funktion\\ \hline

Updates &  über App-Store & Serverseitig & Serverseitig\\ \hline
   				  						 
				
\end{tabular}    
\caption{Veröffentlichung und Installation \cite{PwaNvaWa}}
\label{tab:PwaNvaWaInstallation}
\end{table}


\begin{table}[h]
\centering

\begin{tabular} {|p{3cm}|p{3.5cm}|p{3.5cm}|p{3.5cm}|}
\hline\multirow{3}{*}
 										&PWA  & Native & Web App	\\ \hline
Offline-Zugriff & Verfügbar & Man muss die App einmal online nutzen, dann sollten die Inhalte im Cache offline verfügbar sein. & nicht möglich\\ \hline

Starten im Vollbildmodus & Verfügbar  & Verfügbar & nicht möglich\\ \hline

Kundenbindung &  sehr hoch, Kunden verbringen viel Zeit & App ist wie ein Tap, das macht es für den Kunden leichter zu wechseln & wie \acs{PWA}\\ \hline


				  						 
				
\end{tabular}    
\caption{Zugriff \cite{PwaNvaWa}}
\label{tab:PwaNvaWaZugriff}
\end{table}





Wie in den Tabellen ersichtlich bietet die \acs{PWA} eine Reihe von Vorteilen, z.B. Push Notifikations und Offline-Zugriff, die bei der Benutzung behilflich sein können. Aufgrund dessen bietet sie eine gute Alternative zu den Native App. Probleme machen die Betriebssysteme, da nicht alle Funktionen auf den verschiedenen Systemen zur Verfügung stehen \cite{PwaNvaWa}.
In den nächsten Kapiteln werden die Methoden und APIs in der Theorie und im Kapitel \ref{chap:Implementierung} die praktische Anwendung an einer selbst erstellten App erklärt. 

\section{Web App Manifest}\label{sub:Manifest}
Das App Manifest ist eine JSON Datei die dem Browser verrät, wie sich die \acs{Web-App} bei der Installation auf dem Startbildschirm verhält. Im Manifest werden der Name, der Kurzname, die Größe, das Aussehen der Icons und weitere Eigenschaften definiert. Dessen Zweck ist es der Anwendung auf dem Startbildschirm ihr Aussehen zu verleihen. 
Das App Manifest.json Datei wird in die gleiche Ebene wie die Index.html Datei in das Projekt eingepflegt und über den folgenden Link-Tag im Header implementiert: 

\begin{lstlisting}[language=HTML, caption={Manifest.json} {\cite{Manifest}},label=lst:Manifest.json, xleftmargin=50pt]
<link rel="manifest" href="/<Dateinname>">
\end{lstlisting}

Bei Anwendungen mit mehreren \acs{HTML}-Seiten muss der Link-Tag auf jeder Seite eingefügt werden. 
Im Listening \ref{lst:Manifest.jsonBsp} ist ein Auszug vom Aufbau dargestellt:
	\begin{lstlisting}[language=json, firstnumber=1, caption={Manifest in das Projekt implementieren} {\cite{Manifest}},label=lst:Manifest.jsonBsp, xleftmargin=50pt]
  "name":"PWA Smart Home RMJ",
  "short_name":"PWA_SHL_RMJ",
  "start_url":"./",
  "scope":".",
  "display":"standalone",
  "background_color":"#003399",
  "theme_color":"#3F51C5",
  "icons":[
    {
      "src":"./static/img/light48.png",
      "type":"image/png",
      "sizes":"48x48"
    }
  ]

\end{lstlisting}











%\chapter{Entwurf}
\thispagestyle{standard}
\pagestyle{standard}
\renewcommand{\footrulewidth}{0.4pt}
\lfoot{\small Refik Kerimi}

In diesem Kapitel werden die Muster und die Anforderungen bzw. die Umsetzung der \acl{PWA}en (\acs{PWA}) im Allgemeinen betrachtet.


\section{Übersicht PWA}\label{sub:Übersicht PWA}






\chapter{Implementierung}\label{chap:Implementierung}
\thispagestyle{standard}
\pagestyle{standard}
\lfoot{\small Refik Kerimi}

In diesem Kapitel wird die Umsetzung der Applikation beschrieben.

\section{Anforderungsanalyse}\label{sub:Anforderungsanalyse}
Die zu entwickelnde Smart Home Applikation muss das Verhalten einer \acs{PWA} aufweisen. Das heißt es müssen die Attribute aus Kapitel \ref{chap:ProgressiveWebapplikationen} eingebaut werden und danach getestet werden siehe Kapitel \ref{chap:Funktionstest}. Außerdem sollen auch APIs entwickelt werden die es möglich machen mit Services von Drittanbietern die Temperatur zu regeln, mit Hilfe von Geolocation API den Standort zu ermitteln um das Garagentor über GPS automatisch öffnen zu können. Ebenso soll das Steuern der Beleuchtung möglich sein. Diese Applikation soll Offline so weit es geht verwendbar sein. Die Anwendung muss Plattformunabhängig und responsive sein.

\section{Umsetzung der Anforderungen}\label{sub:Umsetzung der Anforderungen}
Zur Umsetzung der Anforderungen aus Kapitel \ref{sub:Anforderungsanalyse} wurden für  User Interfac ReactJS\footnote{https://reactjs.org/docs/getting-started.html} und als CSS Framework Semantic-UI\footnote{https://react.semantic-ui.com/introduction} verwendet. Sematic-UI soll sicherstellen das die Applikation responsives Verhalten aufweist und für alle Bildschirmgrößen geeignet ist. Um die Daten zu versenden, aufzurufen und zu speichern wurde das JSON Key/Value Format, die Fetch API und der Browser Cache verwendet.
Als Browser diente der Google Chrome Version 67.
Die nicht fertigen Funktionen wurden mit Mockups dargestellt, um einen Eindruck zu vermitteln wie die App in Zukunft aussehen soll.

\section{Ausgewählte Programmiersprache und IDE}
Als Programmiersprache wurde \acl{JS} (\acs{JS}) ausgewählt. 
Als Entwicklungsumgebung wurde Webstorm (Version 2018.2) von Jetbrains verwendet. 
Weitere verwendete Tools und Frameworks wurden in Kapitel \ref{sub:Umsetzung der Anforderungen} beschrieben.
\newpage
\section{Ordnerstruktur}
Die zwei wichtigsten Dateien befinden sich wie in Abbildung \ref{fig:OrdnerStrucktur} zu sehen ist, im app Verzeichnis.
Ebenfalls wichtig ist das app.js File das unter \textit{/app/src/js/app.js} zu finden ist.

\begin{figure}[h]
	\centering
	\includegraphics[width=10cm]{BilderAllgemein/Implementierung/OrdnerStrucktur.jpg}\medskip
	\caption{Ordner Strucktur}
	\label{fig:OrdnerStrucktur}
\end{figure} 


\section{Manifest}
Das Manifest wird im Root Folder eingefügt und ist mit der Endung \acs{JSON} deklariert. Der genau Pfad ist \textit{/app/manifest.json}. Wie in Kapitel \ref{sub:Manifest} schon beschrieben wurde, definiert das Manifest File das Aussehen des Icons am Startbildschirm, z.B. den Einstiegspunkt der App und den Namen. 
Die Listing \ref{lst:Manifest.json} zeigt weitere wichtige Eigenschaften:
\newpage
\begin{lstlisting}[language=json, firstnumber=1, caption={Manifest in das Projekt implementieren} {\cite{Manifest}},label=lst:Manifest.json, xleftmargin=50pt]
{
  {
  "name":"PWA Smart Home",
  "short_name":"PWA_SHL_RMJ",
  "start_url":"./",
  "scope":".",
  "display":"standalone",
  "background_color":"#003399",
  "theme_color":"#3F51C5",
  "description":"Keep running with PWA",
  "dir":"ltr",
  "lang":"de-DE",
  "orientation":"portrait-primary",
  "icons":[
    {
      "src":"./static/img/light48.png",
      "type":"image/png",
      "sizes":"48x48"
    },
    {
      "src":"./static/img/light512.png",
      "type":"image/png",
      "sizes":"512x512"
    }
  ]
}
}
\end{lstlisting}
\clearpage
\section{Add to Homescreen}
Damit der Add to Homescreen Banner erscheint müssen die Bedingungen aus dem Kapitel \ref{sub:AddtoHomescreen}erfüllt werden.
Nach Erfüllung dieser Forderungen wird der beforeinstallprompt Event wie in Listing \ref{lst:AddtoHomescreen} aufgerufen und in die deferredPrompt variable gespeichert.

Verzeichnis: \textit{/app/src/js/app.js}

\begin{lstlisting}[language=JavaScript, firstnumber=1, caption={Add to Homescreen Funktion} ,label=lst:AddtoHomescreen, xleftmargin=50pt]
var deferredPrompt;

window.addEventListener('beforeinstallprompt', event => {
    event.preventDefault();
    deferredPrompt = event;
    return false;
});
%\end{lstlisting}

Diese Callbackfunktion wird in die \textit{app.js} Datei implementiert. 

\section{Service Worker und Cache API}
Das Herzstück der Applikation ist der im Kapitel \ref{chap:Basistechnologien} beschriebene Service Worker. Als erstes muss validiert werden ob der Service Worker vom Browser unterstützt wird und danach kann dieser registriert, installiert und aktiviert werden.
In den folgenden Listings \ref{lst:Registrierung}, \ref{lst:Installation} und \ref{lst:Aktivierung} sind die Funktionen, die für den \acs{SW} von Bedeutung sind aufgeführt und die wichtigsten Teile beschrieben.

Verzeichnis: \textit{/app/src/js/app.js}

\begin{lstlisting}[language=JavaScript, firstnumber=1, caption={Registrierung} ,label=lst:Registrierung, xleftmargin=50pt]
//Registrierung und Validierung vom Service Worker
if ('serviceWorker' in navigator && 'Notification' in window) {
    console.log('Service Worker and Notification is supported')
    navigator.serviceWorker.register('/sw.js')
        .then(reg => {
            console.log('Service worker registered!', reg);

        })
        .catch(err => console.log(err));
}
\end{lstlisting}

\clearpage
Die Registriermethode \textit{register()} bekommt als Parameter die Service Worker Datei mitgegeben.

Verzeichnis: \textit{/app/sw.js}

\begin{lstlisting}[language=JavaScript, firstnumber=1, caption={Installation} ,label=lst:Installation, xleftmargin=50pt]
self.addEventListener('install', event => {
    console.log('[Service Worker] Installing Service Worker ....', event);
    event.waitUntil(
        caches.open(cacheName)
            .then(cache => {
                console.log('[Service Worker] Precaching App Frame');
                cache.addAll(filesToCache);
            })
    )
});
}
\end{lstlisting}

Nach dem erneuten Laden der Anwendung wird im Listing \ref{lst:Installation} der Installations Eventlistener aufgerufen. Dieser installiert den Service Worker und cached die angegebenen Dateien um die App offline verwenden zu können. 

Durch die Funktion \textit{waitUntil()} wird mit der Installation gewartet, bis die Dateien die dieser Funktion als Parameter mitgegeben wurden, gecacht wurden. Durch die Cachemethode \textit{cache.All()} werden alle angegeben Dateien aufgerufen und es müssen nicht alle Dateien einzeln eingeben werden.
Nach der Installation wird der Service Worker aktiviert und kann dann vom Browser verwendet werden.

Verzeichnis: \textit{/app/sw.js}

\begin{lstlisting}[language=JavaScript, firstnumber=1, caption={Aktivierung} ,label=lst:Aktivierung, xleftmargin=50pt]
self.addEventListener('activate', event => {
    console.log('[Service Worker] Activating Service Worker ....', event);
    return self.clients.claim();
})
\end{lstlisting}

Durch die \textit{self.clients.claim()} Methode in Zeile 3 wird sichergestellt, dass der Service Worker nur installiert wird wenn alle Bedingungen erfüllt wurden.

In der Callback-Funktion \textit{respondWith()} werden die Daten aufgerufen die match-Methode überprüft ob die Daten sich im Cache befinden.
Um den Event aufzurufen werden Promises für asynchrone Aufrufe verwendet .
Um Daten aus dem Netzwerk aufzurufen die nicht im Cachespeicher vorhanden sind wird über den fetch Event aufgerufen und überprüft wie in Abbildung \ref{lst:fetch} zu sehen.

Verzeichnis: \textit{/app/sw.js}

\begin{lstlisting}[language=JavaScript, firstnumber=1, caption={fetch Calbackfunktion} ,label=lst:fetch, xleftmargin=50pt]
self.addEventListener('fetch', event => {
    event.respondWith(
        caches.match(event.request)
            .then(response => {
                if (response) {
                    return response;
                }else{
                    return fetch(event.request);
                }
            })
    );
});
\end{lstlisting}


	 
\clearpage
\section{Offline Modus}
Eine der wichtigsten Aufgaben des Caches vom Service Worker ist der Offlinemodus oder das Arbeiten bei schlechter Internetverbindung. Der Cache beinhaltet im Grunde das Index.html File CSS und Bilder oder Icons. Bei der Entwicklung der Anwendung wurden statische Dateien wie im Listing \ref{lst:Cache} verwendet. Die benötigten Dateien wurden in die \textit{let fileToCache} Variable gespeichert und im Listing \ref{lst:Installation} und der \textit{caches.addAll()}-Methode als Parameter mitgegeben.


\begin{lstlisting}[language=JavaScript, firstnumber=1, caption={Cache} ,label=lst:Cache, xleftmargin=50pt]
let cacheName = 'precache';
let filesToCache = [
    '/',
    '/index.html',
    '/src/js/app.js',
    '/static/img/light48.png',
    '/static/img/dashboard-mockup.jpg',
    '/static/img/bulp.jpeg',
    '/static/img/garage.jpeg',
    '/static/img/Graph_Heizung.JPG',
    '/manifest.json'
];
\end{lstlisting}


\begin{figure}[H]
	\centering
	\includegraphics[width=16cm]{BilderAllgemein/Implementierung/Cache.jpg}\medskip
	\caption{Cache}
	\label{fig:Cache}
\end{figure}  

\begin{figure}[H]
	\centering
	\includegraphics[width=8cm]{BilderAllgemein/Implementierung/aufruf_SW_Browser.jpg}\medskip
	\caption{Datenaufruf Service Worker}
	\label{fig:aufrufSWBrowser}
\end{figure}  

In der Abbildung \ref{fig:aufrufSWBrowser} und \ref{fig:Cache} sind die gecachten Files vom Service Worker im Netzwerk und im Cache zu erkennen.



\clearpage
\section{Push Notifications}
Die Push Benachrichtigungen dienen dem Nutzer dazu, bei Änderungen im Haushalt benachrichtigt zu werden, z.B.: wenn das Licht defekt ist oder die Verbindung zur Heizung ausfällt.

Die Benachrichtigungen wurden in diesem Projekt nicht vom Server aus aufgerufen sondern sind über folgendes Listing im Pfad: \textit{/app/src/app.js} eingetragen. 
Nach der Überprüfung im Listing \ref{lst:Registrierung} ob die Pushfunktion im Browserscope vorhanden ist, wird in Listing \ref{lst:notificationsListener} der  Eventlistener über eine Callbackfunktion aufgerufen. Bei diesem Prototypen sind die Benachrichtigungen im JSON-File\\ (Verzeichnis: \textit{/app/src/data/notification.json}) und werden über eine fetch-API aufgerufen. Durch die \textit{openWindow()}-Funktion in Zeile 8 wird die Seite mit dem Benachrichtigungstext geöffnet.  
\\
\begin{lstlisting}[language=JavaScript, caption={Geolocation  Support} {\cite{PushNotifikationServer}},label=lst:notificationsListener, xleftmargin=50pt]
self.addEventListener('notificationclick', e => {
    let notification = e.notification;
    let action = e.action;

    if (action === 'close') {
        notification.close();
    } else {
        clients.openWindow('http://localhost:8080/benachrichtigungen');
        e.notification.close()

    }
});
\end{lstlisting}
\clearpage

\section{Geolocation API}
Um zu zeigen, dass der Zugriff auf die Geräte-APIs möglich ist wurde in der Applikation eine Funktion hinzugefügt, die den genauen Standpunkt über die Geolocation-API ermittelt. Diese Funktion könnte als Garagenöffner oder zum Einschalten der Heizung nützlich sein. Wenn sich der Nutzer dem Haus nähert könnte über die GPS Daten das Tor geöffnet werden ohne, dass der Benutzer über eine HCI eingreifen muss. 
Als erstes wird im Listing \ref{lst:GeolocationSupport} der Support des Browsers überprüft.\\
Verzeichnis: \textit{/app/sw.js}

\begin{lstlisting}[language=JavaScript, caption={Geolocation  Support} {\cite{UserLocation}},label=lst:GeolocationSupport, xleftmargin=50pt]
if (navigator.geolocation) {
  console.log('Geolocation is supported!');
}
else {
  console.log('Geolocation is not supported for this Browser/OS.');
}
\end{lstlisting}

In der Abbildung \ref{fig:GeoSupport} sieht man, dass der Chrome Browser diese Geolocation Funktion unterstützt und man sieht auch die genauen Koordinaten.

\begin{figure}[h]
	\centering
	\includegraphics[width=14cm]{BilderAllgemein/Geolocation/GeoSupport.jpg}\medskip
	\caption{Konsolenmeldung Geolocation}
	\label{fig:GeoSupport}
\end{figure}








\chapter{Funktionstest/Validierung}\label{chap:Funktionstest}
\thispagestyle{standard}
\pagestyle{standard}
\lfoot{\small Refik Kerimi}

\section{Ausgangsbedingung und Ausgrenzung}
Die \acs{PWA} wird die Native Apps wahrscheinlich nie ersetzten, aber sie haben sicher das Potenzial eine feste Strategie bei der Entwicklung von \acs{Web-App}s zu werden.




\chapter{Zusammenfassung und Ausblick}
\thispagestyle{standard}
\pagestyle{standard}
\renewcommand{\footrulewidth}{0.4pt}
\lfoot{\small Refik Kerimi}


%\chapter{Einleitung}

\thispagestyle{standard}
\pagestyle{standard}
\section{Motivation}

Die zunehmende Vernetzung und die immer stärkere Integration von Systemen macht
auch vor dem Bereich der Automatisierungstechnik nicht Halt. So ist hier in den letzten
Jahren ein steigender Bedarf an Interkonnektivität der einzelnen, vormals autonomen
Systeme zu verzeichnen. Vorschub leistet dieser Bewegung, dass durch die gestiegene
Verfügbarkeit von breitbandigen Verbindungen auch der Wunsch nach höherer Integrationsdichte
gewachsen ist. Immobilien, Ladengeschäfte, Produktionsstätten oder ganze
Firmenstandorte werden zunehmend als Verwaltungseinheiten gesehen, für die eine Administration
und Überwachung aus der Ferne oder die Einbindung in Firmennetze erfolgen
soll.
Auch der immer häufigereWunsch, Geschäftsprozesse automatisiert in Software abzubilden,
lässt es sinnvoll erscheinen, auf verschiedene Systeme eines Standorts aus der Ferne
als Entität zugreifen zu können. Damit verändern sich auch die Anforderungen an die
Schnittstellen eines Systems. Waren früher systemindividuelle Bediener vor Ort für die
Betreuung ausreichend, so soll dies heute immer öfter auch aus der Ferne zentral über
eine standardisierte Schnittstelle möglich sein. Somit hält die Problematik der verteilten
Anwendungen auch in vormals oft vollständig unabhängigen Systemen Einzug.
Gerade in der Automatisierungstechnik macht sich bemerkbar, dass der scheinbare Wildwuchs
an konkurrierenden Bus Systemen der letzten Jahre zu einer extremen Inkompatibilität
der Systeme geführt hat. Systeme die nun wieder eine Einheit im Sinne der remoten
Verwaltung bilden sollen. An Hand der Energieverwaltung eines Gebäudes oder
einer Produktionsstätte lässt sich dies veranschaulichen. Für eine zentrale Überwachung



\section{Zielsetzung}

\section{Anmerkung}

\chapter{Grundlagen}

\section{OpenNES}

\subsection{Architektur}

\subsection{}

\section{ModBus Protokoll}

\section{Abbildungen}

In Moby-Dick geht es in erster Linie um die Jagd auf einen weißen Pottwal \footnote{nach: https://de.wikipedia.org/wiki/Moby-Dick}. Kurze Zitate (unter drei Zeilen) müssen mit Anführungszeichen gekennzeichnet werden. Außerdem müssen die Quelle sowie die Seite angegeben werden. Ein Beispiel für ein kurzes Zitat: \glqq Komisch. Manch einer von uns wünschte sich, er lebe auf einer Südseeinsel.\grqq \cite{MELVILLE:MOBYDICK1997} (Seite 100).

  \begin{quote}
"Das Buch hier Lieblingsbuch. Viele Blätter viele, schöne Bilder. Du kennen Worte?"\newline
"Ja"\newline
"Ich kennen Bilder. Das ein Wal. Du lesen Worte!"\newline
"Durch das Herz des Wals strömt mehr Flüssigkeit als durch das große Wasserleitungsrohr unter der London Bridge, jedoch strömt das Wasser nicht so stark, wie das Blut, das vom Herz des Wals pocht."\newline
"Du gut, ich danken dir."  \upshape \cite{MELVILLE:MOBYDICK1997} (Seite 500)
  \end{quote}

Ein sehr praktisches Package ist cleveref. Es automatisiert und erleichtert das setzen von Referenzen ungemein. Als Beispiel wird eine Referenz auf das FHS-Logo gesetzt siehe \cref{FIG_LOGO}.

Für das Verfassen von wissenschaftlichen Arbeiten können eine Vielzahl an Quellentypen herangezogen werden. Beispiele hierfür sind Leitfäden (Manuals) \citep{RFC2828} \citep{80211i} \citep{80211} \cite{X800} \citep{TR102377} \citep{EN301893} \citep{PUB197} \citep{PUB74}, Bücher \citep{Fis04a} \citep{Rei05a} \citep{Tan00a} \citep{Ste04a} \citep{GMS00a} \citep{HL98a}, Sammelbände \citep{EHL00a} \citep{Sch94a}, Journal-Artikel \citep{TM03a} \citep{CP03a}, Konferenz-Proceedings \citep{HCB00a} \citep{KBW04a} \citep{KSW04a} \citep{HK05a}, Internetmagazine \citep{Eke05a}, Webquellen \cite{nist} \cite{php} \cite{BDKMT93a} \cite{IDSSM} sowie Diplomarbeiten und Dissertationen \citep{Sch98} \cite{Hae94a}.

Als Beispiel für eine Abkürzung wird hier \ac{ADF} angeführt. Das Package schreibt automatisch das erste Vorkommen der Abkürzung aus. Die zweite Verwendung von \ac{ADF} wird also abgekürzt. Ist ein Ausschreiben einer Abkürzung gewünscht wird der acl-Befehl verwendet. Dies führt zu \acl{ADF}. Abkürzungen müssen in der Datei \glqq 05Abkuerzungsverzeichnis\grqq angegeben werden.

\section{Quelltext}

\texttt{printf("Hallo Welt")} für Ausschnitte von Sourcecode innerhalb von Text

\lstset{escapeinside={\%*}{*)},numbers=none}%oder numbers=left
\begin{lstlisting}[language=C,
caption=Beispiel-Listing,
label=LST_SAMPLE]
serverTCP = new TcpListener(IPAddress.Parse(serverIP), serverPort);
\end{lstlisting}

\lstset{escapeinside={\%*}{*)},numbers=left}
\lstinputlisting[language=Matlab, caption=Einfaches Matlabprogramm in einer Datei, label=list:hello.m]{Listings/hello.m}
\section{Text}

\section{Refik}

\section{Bilder}

\begin{figure}[H]
\begin{center}
	\includegraphics[scale=0.4]{BilderAllgemein/Logo.jpg}
\end{center}
	%\includegraphics[width=\textwidth]
	%\end{center}
	% Title
	\caption{Das FHS-Logo}
	% Unique name: identifier for referencing
	\label{FIG_LOGO}
\end{figure}

\section{Formeln}

Formeln sind für jeden Abschnitt rechtsbündig von dieser zu nummerieren, um einen späteren Bezug in der Arbeit zu gewährleisten. Formeln werden üblicherweise in "`Computer Modern Roman"' (\LaTeX{}-Standard) gesetzt. In diesem Template wird die Formel-Schrift bzw. das Package \texttt{eulervm} verwendet. Abgesetzte Formeln werden in \LaTeX{} durch die 
\emph{equation} Umgebung definiert. Formelausdrücke innerhalb von Textabschnitten erhält man durch \$Formel\$.

\subsection*{Beispiel}
%
Der \emph{Sinus cardinalis} oder sinc-Funktion ist eine mathematische Funktion $f$, welche in nicht-normierter Version als

\begin{equation}
	f(x) := \frac{\sin(x)}{x}
	\label{eq:bsp}
\end{equation}

definiert wird. In der digitalen Signalverarbeitung findet meistens nachfolgende normierte Version $\mathrm{si}(x)$ oder $\mathrm{sinc}(x)$ Anwendung \cite{x1}, \cite{x2}. Für eine Visualisierung dieser Funktionen siehe Abb.~\ref{FIG_LOGO}.

\begin{equation}
	f(x) := \frac{\sin(\pi x)}{\pi x}
	\label{EQ_SAMPLE}
\end{equation}

\section{Beispiel für Tabellen}
%
Es empfiehlt sich, für Tabellen die Standard-\LaTeX{}-Umgebung \emph{tabular} zu verwenden. Bei Bedarf können natürlich auch Erweiterungen (z.B.~\emph{tabularx} oder \emph{array}) zur Anwendung kommen. Eine mögliche Darstellung zeigt Tabelle \ref{Table_Sinc}.

\begin{table}[h!]%
	\begin{center}
	
		\begin{tabular}{|r|r|r|}
			\firsthline
			$x$&$\mathrm{sinc}(x)$&$\mathrm{sin}(x)$\\\hline\hline
			$-0.5$&0.6366&-0.4794\\\hline
			$0$&1.0000&0\\\hline
			$0.5$&0.6366&0.4794\\\hline
		\end{tabular}
		\caption{Zwei Werte der Sinc-Funktion}
		\label{Table_Sinc}
	\end{center}
\end{table}













%%%%%%%%%%%%%%%%%%%%%%%%%%%%%%%%%%%%%%%%%%%%%%%%%%%%%%%%%%%%%%%%%%%%%%%%%%%%%%%%%%%%%%%%%%%%%%%%%%%%%%%%%%%%
% LITERATURVERZEICHNIS

\interlinepenalty = 10000 % Literatureinträge: Absätze zusammenhalten
\clearpage
\addcontentsline{toc}{chapter}{Literatur}
\singlespace
\raggedright
\bibliography{12Literatur}
\interlinepenalty = 100

%%%%%%%%%%%%%%%%%%%%%%%%%%%%%%%%%%%%%%%%%%%%%%%%%%%%%%%%%%%%%%%%%%%%%%%%%%%%%%%%%%%%%%%%%%%%%%%%%%%%%%%%%%%%
% ANHÄNGE

\begin{appendix}
%\include{Anhang_Mathematik}                 % A
%\include{Anhang_FormatDerParameterdateien}  % B
%\include{Anhang_Quelltexte}                 % C
%\include{Anhang_Datenblaetter}              % D
%\include{Anhang_Glossar}                    % E
\end{appendix}

\end{document}